% Shared preamble of intro.tex (the text) and exercises.tex (exercises and solutions).
\usepackage[T1]{fontenc}
\usepackage[utf8]{inputenc}
% one typeface for everything: Palatino for text and mathematics, with real small caps
\usepackage[sc]{mathpazo}
\linespread{1.08}
\usepackage{amsmath,amssymb,amsthm}
\usepackage{microtype}
% a readable measure: about 70 characters a line
\usepackage[a4paper,textwidth=140mm,top=32mm,bottom=34mm,headsep=10mm]{geometry}
\usepackage{xcolor}
\definecolor{linkblue}{rgb}{0.10,0.20,0.45}
\usepackage[colorlinks,linkcolor=linkblue,citecolor=linkblue,urlcolor=linkblue,filecolor=linkblue]{hyperref}
% quiet running heads: the current section on the left, the page on the right, in small caps; no rules
\usepackage{fancyhdr}
\pagestyle{fancy}
\fancyhf{}
\makeatletter\newcommand{\topsection}{\expandafter\@leftmark\topmark\@empty\@empty}\makeatother  % the section in effect at the top of the page
\fancyhead[L]{\small\scshape\nouppercase{\topsection}}
\fancyhead[R]{\small\thepage}
\renewcommand{\headrulewidth}{0pt}
\renewcommand{\sectionmark}[1]{\markboth{#1}{}}
\setlength{\parindent}{1.2em}

\newcommand{\op}{\cdot}
\newcommand{\cop}{\circ}
\newcommand{\N}{\mathbb{N}}
\newcommand{\Z}{\mathbb{Z}}
\newcommand{\proj}{\textsc{duet}}
\newcommand{\file}[1]{\mbox{\textit{#1}}}  % file names: italic, same typeface, never hyphenated

% headings: only two sizes (large for sections, normal for run-in paragraphs), all bold, same typeface
\makeatletter
\renewcommand\section{\@startsection{section}{1}{\z@}{-3.2ex \@plus -1ex \@minus -.2ex}{1.4ex \@plus .2ex}%
  {\normalfont\large\bfseries}}
\renewcommand\paragraph{\@startsection{paragraph}{4}{\z@}{1.8ex \@plus .5ex \@minus .2ex}{-0.6em}%
  {\normalfont\normalsize\bfseries}}
% a part is a quiet small-caps line, not a page-sized heading
\newcommand{\partline}[2]{\par\bigskip\addcontentsline{toc}{part}{Part #1\quad #2}%
  {\centering\normalfont\large\scshape Part #1\quad #2\par}\nobreak\medskip}
% table of contents: parts in small caps at normal size, sections without extra space between them
\renewcommand*\l@part[2]{\addpenalty{-\@highpenalty}\vskip .8em \@plus\p@
  {\parindent\z@\rightskip\@pnumwidth\parfillskip-\@pnumwidth\normalfont\scshape #1\hfil\hb@xt@\@pnumwidth{\hss #2}\par}}
\let\intro@lsection\l@section
\renewcommand{\l@section}[2]{\begingroup\let\addvspace\@gobble\intro@lsection{#1}{#2}\endgroup}
% lists a little tighter than the default
\def\@listi{\leftmargin\leftmargini \parsep\z@ \topsep 5\p@ \@plus 2\p@ \@minus 1\p@ \itemsep 3\p@ \@plus 1\p@}
\let\@listI\@listi
\makeatother

% every statement upright, with a bold label
\newtheoremstyle{uniform}{\medskipamount}{\medskipamount}{\normalfont}{}{\bfseries}{.}{0.5em}{}
\theoremstyle{uniform}
\newtheorem{theorem}{Theorem}[section]
\newtheorem{proposition}[theorem]{Proposition}
\newtheorem{definition}[theorem]{Definition}
\newtheorem{example}[theorem]{Example}
\newtheorem{exercise}{Exercise}

